% GENERATED FILE. Edit canonical lessons, then run npm run generate:books.
% canonical-source-hash: fnv1a64:89c5449f882dc721
% canonical-lessons: IT-C338-frusta, IT-C338-sbattere, IT-C338-montare, IT-C338-amalgamare, IT-C338-mestolo, IT-R338-ripasso-337-338, IT-R338-ripasso-final-338

\chapter{Tutto insieme}
\label{ch:it-tutto-insieme}
\hlchaptermodality{338}

\begin{chapteropening}
\textbf{By the end of this chapter:} \emph{I can describe mixing: the whisk, beating eggs, whipping, blending and the ladle.}

Complete the last lesson of chapter 338: I can describe mixing: the whisk, beating eggs, whipping, blending and the ladle.
\end{chapteropening}

\section[la frusta]{la frusta — the whisk; the whip}

\label{lesson:IT-C338-frusta}



\begin{quote}
Before the new one: say the Italian for to stuff, then the Italian for to garnish.
\end{quote}

\subsection*{You'll want to know: la frusta}
\textbf{la frusta} — \textquotedblleft{}the whisk; the whip\textquotedblright{}.

\begin{quote}
Prendi la frusta, non il cucchiaio.
\end{quote}

\begin{quote}
\emph{Take the whisk, not the spoon.}
\end{quote}

\emph{La frusta elettrica} — the electric whisk.

\begin{cousinweb}
\emph{frusta} is traced in EtymDB to Latin \emph{fūstis} \textquotedblleft{}a stick, a club\textquotedblright{}.
\end{cousinweb}

\begin{grammarlens}[title={What you've built}]
One more word for telling, explaining and describing at B1.
\end{grammarlens}

\subsection*{Your turn}
\begin{itemize}
\raggedright
  \item \emph{Say it:} \emph{la frusta}
  \item \emph{Say it:} \emph{la frusta}, once more
  \item \emph{Say it:} a sentence of your own with \emph{la frusta}
  \item \emph{Recall:} say the Italian for to stuff, then the Italian for to garnish, then say \emph{la frusta} again
\end{itemize}

\begin{tcolorbox}[breakable,colback=teal!4,colframe=teal!35!black,title={Before you move on}]
What does la frusta mean? (\textquotedblleft{}The whisk; the whip\textquotedblright{}.) Say it once more.
\end{tcolorbox}
\section[sbattere]{sbattere — to beat (eggs); to bang, to slam}

\label{lesson:IT-C338-sbattere}



\begin{quote}
Before the new one: say the Italian for to garnish, then the Italian for the whisk.
\end{quote}

\subsection*{You'll want to know: sbattere}
\textbf{sbattere} — \textquotedblleft{}to beat (eggs); to bang, to slam\textquotedblright{}.

\begin{quote}
Sbatti le uova con lo zucchero.
\end{quote}

\begin{quote}
\emph{Beat the eggs with the sugar.}
\end{quote}

Also \emph{sbattere la porta} — to slam the door.

\begin{cousinweb}
\emph{sbattere} is \emph{s-} with \emph{battere} \textquotedblleft{}to beat\textquotedblright{}, from Latin \emph{battuere} (Lewis \& Short).
\end{cousinweb}

\begin{grammarlens}[title={What you've built}]
One more word for telling, explaining and describing at B1.
\end{grammarlens}

\subsection*{Your turn}
\begin{itemize}
\raggedright
  \item \emph{Say it:} \emph{sbattere}
  \item \emph{Say it:} \emph{sbattere}, once more
  \item \emph{Say it:} a sentence of your own with \emph{sbattere}
  \item \emph{Recall:} say the Italian for to garnish, then the Italian for the whisk, then say \emph{sbattere} again
\end{itemize}

\begin{tcolorbox}[breakable,colback=teal!4,colframe=teal!35!black,title={Before you move on}]
What does sbattere mean? (\textquotedblleft{}To beat (eggs); to bang, to slam\textquotedblright{}.) Say it once more.
\end{tcolorbox}
\section[montare]{montare — to whip (cream, eggs); to climb; to assemble}

\label{lesson:IT-C338-montare}



\begin{quote}
Before the new one: say the Italian for the whisk, then the Italian for to beat.
\end{quote}

\subsection*{You'll want to know: montare}
\textbf{montare} — \textquotedblleft{}to whip (cream, eggs); to climb; to assemble\textquotedblright{}.

\begin{quote}
Monta le uova con la frusta per dieci minuti.
\end{quote}

\begin{quote}
\emph{Whip the eggs with the whisk for ten minutes.}
\end{quote}

\emph{Montare a neve} — to beat until stiff. Also \emph{montare un mobile} — to put furniture together.

\begin{cousinweb}
\emph{montare} is built on \emph{monte} \textquotedblleft{}mountain\textquotedblright{}, from Latin \emph{mōns, montis} (Lewis \& Short): to go up, and in the kitchen, to make it rise.
\end{cousinweb}

\begin{grammarlens}[title={What you've built}]
One more word for telling, explaining and describing at B1.
\end{grammarlens}

\subsection*{Your turn}
\begin{itemize}
\raggedright
  \item \emph{Say it:} \emph{montare}
  \item \emph{Say it:} \emph{montare}, once more
  \item \emph{Say it:} a sentence of your own with \emph{montare}
  \item \emph{Recall:} say the Italian for the whisk, then the Italian for to beat, then say \emph{montare} again
\end{itemize}

\begin{tcolorbox}[breakable,colback=teal!4,colframe=teal!35!black,title={Before you move on}]
What does montare mean? (\textquotedblleft{}To whip (cream, eggs); to climb; to assemble\textquotedblright{}.) Say it once more.
\end{tcolorbox}
\section[amalgamare]{amalgamare — to blend, to mix until smooth}

\label{lesson:IT-C338-amalgamare}



\begin{quote}
Before the new one: say the Italian for to beat, then the Italian for to whip.
\end{quote}

\subsection*{You'll want to know: amalgamare}
\textbf{amalgamare} — \textquotedblleft{}to blend, to mix until smooth\textquotedblright{}.

\begin{quote}
Amalgama bene tutti gli ingredienti.
\end{quote}

\begin{quote}
\emph{Blend all the ingredients well.}
\end{quote}

A recipe word: until the mixture is even.

\begin{cousinweb}
\emph{amalgamare} is built on \emph{amalgama}, from Medieval Latin \emph{amalgama}, a word of the alchemists; its deeper origin is not settled (EtymDB has no entry).
\end{cousinweb}

\begin{grammarlens}[title={What you've built}]
One more word for telling, explaining and describing at B1.
\end{grammarlens}

\subsection*{Your turn}
\begin{itemize}
\raggedright
  \item \emph{Say it:} \emph{amalgamare}
  \item \emph{Say it:} \emph{amalgamare}, once more
  \item \emph{Say it:} a sentence of your own with \emph{amalgamare}
  \item \emph{Recall:} say the Italian for to beat, then the Italian for to whip, then say \emph{amalgamare} again
\end{itemize}

\begin{tcolorbox}[breakable,colback=teal!4,colframe=teal!35!black,title={Before you move on}]
What does amalgamare mean? (\textquotedblleft{}To blend, to mix until smooth\textquotedblright{}.) Say it once more.
\end{tcolorbox}
\section[il mestolo]{il mestolo — the ladle}

\label{lesson:IT-C338-mestolo}



\begin{quote}
Before the new one: say the Italian for to whip, then the Italian for to blend.
\end{quote}

\subsection*{You'll want to know: il mestolo}
\textbf{il mestolo} — \textquotedblleft{}the ladle\textquotedblright{}.

\begin{quote}
Aggiungi l'acqua un mestolo alla volta.
\end{quote}

\begin{quote}
\emph{Add the water one ladleful at a time.}
\end{quote}

Wooden or steel; \emph{un mestolo di brodo} is a ladleful of broth.

\begin{cousinweb}
\emph{mestolo} is built on \emph{mestare} \textquotedblleft{}to stir\textquotedblright{} (EtymDB), from Vulgar Latin \emph{mixtāre}, built on \emph{miscēre} \textquotedblleft{}to mix\textquotedblright{}.
\end{cousinweb}

\begin{grammarlens}[title={What you've built}]
One more word for telling, explaining and describing at B1.
\end{grammarlens}

\subsection*{Your turn}
\begin{itemize}
\raggedright
  \item \emph{Say it:} \emph{il mestolo}
  \item \emph{Say it:} \emph{il mestolo}, once more
  \item \emph{Say it:} a sentence of your own with \emph{il mestolo}
  \item \emph{Recall:} say the Italian for to whip, then the Italian for to blend, then say \emph{il mestolo} again
\end{itemize}

\begin{tcolorbox}[breakable,colback=teal!4,colframe=teal!35!black,title={Before you move on}]
What does il mestolo mean? (\textquotedblleft{}The ladle\textquotedblright{}.) Say it once more.
\end{tcolorbox}
\section[(dialogue)]{Ripasso: passo dopo passo}

\label{lesson:IT-R338-ripasso-337-338}



\begin{quote}
Say the Italian for to dress and to salt.
\end{quote}

\subsection*{Your turn}
Say these aloud:

\begin{itemize}
\raggedright
  \item to dress, to salt, to flavour, to stuff, to garnish
  \item the whisk, to beat, to whip, to blend, the ladle
\end{itemize}

\begin{tcolorbox}[breakable,colback=teal!4,colframe=teal!35!black,title={Before you move on}]
To dress? (\textbf{condire}.) To salt? (\textbf{salare}.) To flavour? (\textbf{insaporire}.) To stuff? (\textbf{farcire}.) To garnish? (\textbf{guarnire}.) The whisk? (\textbf{la frusta}.) To beat? (\textbf{sbattere}.) To whip? (\textbf{montare}.) To blend? (\textbf{amalgamare}.) The ladle? (\textbf{il mestolo}.)
\end{tcolorbox}
\section[(dialogue)]{Tutto insieme — once more}

\label{lesson:IT-R338-ripasso-final-338}



\begin{quote}
Say the Italian for the whisk and to beat.
\end{quote}

\subsection*{Your turn}
\begin{itemize}
\raggedright
  \item \emph{Say it:} the whisk, to beat, to whip, to blend, the ladle
\end{itemize}

\begin{tcolorbox}[breakable,colback=teal!4,colframe=teal!35!black,title={Before you move on}]
The whisk? (\textbf{la frusta}.) To beat? (\textbf{sbattere}.) To whip? (\textbf{montare}.) To blend? (\textbf{amalgamare}.) The ladle? (\textbf{il mestolo}.)
\end{tcolorbox}
